\begin{table}[!t]
\caption{The four-part DELB theorem package and the distinct question
answered by each component. The contribution is their joint organization
for non-anticipating integer prediction; the underlying maximum-entropy and
information-projection ingredients remain prior art.}
\label{tab:delb_theorem_package}
\centering
\footnotesize
\setlength{\tabcolsep}{4pt}
\begin{tabularx}{\textwidth}{
    >{\RaggedRight\arraybackslash}p{0.19\textwidth}
    >{\RaggedRight\arraybackslash}p{0.28\textwidth}
    >{\RaggedRight\arraybackslash}X}
\toprule
\textbf{Component} & \textbf{Mathematical object} &
\textbf{Prediction question answered}\\
\midrule
Exact inverse-lattice floor
&
\(\mathcal H_{\mathbb Z}^{-1}
 [H(X_t\mid\mathcal F_{t-m})]\)
&
What MSE floor is imposed by the target uncertainty and the integer lattice,
independently of the forecasting algorithm?\\

Predictor-dependent refinement
&
\(\mathcal H_{\mathbb Z}^{-1}
 [H(X_t\mid\mathcal F_{t-m})+I(E_t;\mathcal F_{t-m})]\)
&
How is the lower bound strengthened for a particular predictor whose error
still retains information about the available history?\\

Necessary and sufficient attainability conditions
&
\(P_E=q_D\) and \(I(E;\mathcal F)=0\)
&
When does a predictor attain the universal DELB, rather than merely satisfy
the inequality?\\

Two-component prediction slack
&
\(I(E;\mathcal F)+D_{\mathrm{KL},2}(P_E\Vert q_D)\)
&
Is nonattainment attributable to residual history dependence, error-shape
mismatch, or both?\\
\bottomrule
\end{tabularx}
\end{table}
